\section{Discussion}\label{sec:discussion}\label{sec:discussion:scope}
\label{sec:discussion:coefficients}\label{sec:discussion:refusal}\label{sec:discussion:envelope}
The two prescribed-bath experiments show how a common particle
formulation connects external-liquid evaporation, internal-core recession
and residual-component transport. Their signed component and energy
exchanges can supply a future tray calculation. Prescribing the bath
isolates this particle response, while leaving countercurrent vapour,
wall, residence and discharge balances to the equipment model. We first
interpret the Faner comparison, then its meaning for DT practice and what
is needed next.

The calculated Faner particle still dries faster than the measured
charge. Geometry is particularly consequential: image-analysis
particle counts cannot be used directly as dry-mass weights, and a volume-equivalent sphere has a
different surface area from an irregular flake of equal volume. For a
number distribution $p_n(d)$ and common density, the corresponding
mass distribution is proportional to $d^3p_n(d)$. A lognormal family
constrained by the reported mean, range and 1--3-mm fraction provides
a reproducible sensitivity construction, but does not recover the
unreported histogram. Its larger-particle mass fraction can prolong
drying even when small particles heat earlier. The direction and size
of the resulting temperature change must come from class-resolved
integration and the actual thermocouple sampling, rather than from
weights selected to obtain a desired displacement of the curves.

Faner reports preheating, so the observed lack of an initial
warming stage has an experimental basis. The exact initial temperature
and measurement trigger are not specified. A slightly subcooled start
would be a different initial-value experiment, with sensible
heating included in the balances; it cannot be represented by shifting
the calculated temperature alone. The single-particle volume mean also
does not determine the temperature of a holder probe. Simultaneous
loading, spatial temperature and size-resolved observations would help
separate startup, heat exposure and internal transport effects.

For DT practice, the three regimes divide the heat demand.
Indirect heating in PD removes external free hexane before live-steam
contact; steam condensation then supplies heat for rapid solvent removal,
and the remaining retained hexane determines the later approach to the
outlet residual. Industrial DT residence and heat supply must serve these
successive needs \citep{kemper2013solvent}. At the same time, the moist heat
treatment must reduce antinutritional activity while preserving useful
protein quality \citep{witte_nd_soybeanmeal,mustakas1981critical}.
The calculated temperature and moisture histories describe that exposure;
a meal-quality kinetic model is outside the present scope.

The moisture result also matters for energy design: water retained
at DT exit contributes to the load on the downstream dryer, which takes
meal towards its storage moisture \citep{witte_nd_soybeanmeal}.
More condensation can accelerate heat supply and solvent removal, while
more retained water can increase later drying duty. A steam-saving claim
therefore requires the full balance of steam, indirect heat and drying
energy, rather than outlet solvent or moisture alone. These prescribed-bath
experiments expose that coupling without quantifying a plant steam saving.

Two independent equilibrium checks give further context for DT practice.
At literature-motivated last-tray compositions, the dry-shell hexane equilibrium floor is
4.5--66.0 ppm without the separate oil-solution sensitivity and
8.1--112.4 ppm with it (Sec.~\ref*{S-sup:desolventizerfloor}). These
equilibrium quantities do not set the residence needed to reach them.
For retained water the diagnostic relation
\begin{equation}
\aw(\Xw)P_{\mathrm{sat}}^{\mathrm w}(T^*)=P
\label{eq:sorbedbp}
\end{equation}
gives a sorption-elevated boiling range of 100.5--\SI{107.6}{\celsius}.
It combines moisture and temperature observations from different sources;
at the moisture of the temperature study, its prediction is 2.1 K above
the observation (Sec.~\ref*{S-sup:envelopescoring}). Neither this check
nor the calculated DT exit moisture is a validation of the
coupled transient particle or the industrial tower.

In the DT history, condensate wets accessible hexane-depleted material, and a
finite pathway transfers it into the existing retained-water inventory.
Experiments by \citet{jovanovich2003wateruptake} distinguish liquid-water
imbibition from vapour sorption and show that it can include capillarity,
matrix hydration and loosely held interstitial water. Their measurements concern
protein isolates at \SI{20}{\celsius}, not extracted meal in a hot DT.
They motivate an effective liquid-connected pathway but do not supply
our coefficient, justify its 1000-fold mobility enhancement, or establish
that all absorbed liquid is molecularly bound. The present closure keeps
one capacity-bounded retained-water inventory and does not resolve those
microscopic partitions.

The imbibition closure also omits water-induced pore displacement, a
hydraulic saturation front, swelling and extra resistance to hexane escape.
In particular, the external water film is assumed permeable to hexane
without an additional transport resistance
(Section~\ref{sec:model:waterfilm}). Because the film persists for a
substantial time in the finite-contact case, this assumption has
experimental consequences. Independent transient uptake measurements on
the actual extracted meal, with external liquid quantified separately and
several particle sizes, are needed to identify the contact and internal
resistances. Closer agreement with an outlet moisture range alone would
not identify them.

\label{sec:discussion:extrapolation}\label{sec:discussion:geometry}\label{sec:discussion:limits}
Material and numerical limits set the present scope.
The analytic water isotherm
is extrapolated below its original 0.023 kg/kg anchor. Hexane retention
and binary gas cross properties also have temperature and composition
limits, and the transport coefficients are not yet identified from
simultaneous water--hexane measurements. Radius and porosity are fixed. The isobaric
pore closure omits a total-vapour Darcy/Knudsen resistance and assumes
drainage of an immobile connected liquid-hexane core. Numerically, the
two-cell DT and four-cell Faner histories provide complete discrete
experiments, but operator verification and selected timestep checks do not
yet establish whole-history continuum convergence.

\section{Conclusions}\label{sec:conclusions}\label{sec:discussion:measurements}
One conservative particle formulation completes both a 240-s
water-bearing particle experiment in pure superheated hexane and a
1766.5-s prescribed DT history. Component and energy balances
continue through external-liquid depletion, internal-front recession and
core extinction. External condensate is stored separately and exchanges
water and energy with the particle and the prescribed gas.

The declared Faner particle falls from 0.20 to 0.10 kg hexane/kg dry meal
in \SI{\FanerFallingInterval}{\second}, compared with 25.4 s measured. The DT calculation includes liquid-water imbibition
and reaches 18.2\% wet-basis moisture, 111 ppm wet-basis hexane and a
volume-mean particle temperature of \SI{105.0}{\celsius}. Its liquid
contact conductance is derived from one same-material uptake time and its
liquid-water mobility is a declared literature limit; neither is measured
on DT meal.

The model's signed water, hexane and energy fluxes can be coupled to tray
balances for DT design; in that coupled form the formulation is the
particle-level basis for a DT digital twin that can shorten process
design and support model-based control. The results connect the model's
mechanisms to observable trajectories and show which independent
geometry, uptake and temperature measurements are needed before it is
used for predictive equipment calculations.
